\documentclass[12pt,a4paper]{article}
\usepackage[left=1.5cm,right=1.5cm,top=2cm,bottom=2.5cm]{geometry}
\usepackage{multicol}
\usepackage{amsmath,amssymb,amsfonts,latexsym,mathtools}
\usepackage[spanish,es-tabla]{babel}
\usepackage[utf8]{inputenc}
\usepackage[T1]{fontenc}
\usepackage{answers}

% Configuración del paquete answers
\Newassociation{sol}{Solution}{ans}
\newtheorem{ex}{Ejercicio}

\begin{document}

% Encabezado Institucional
\begin{center}
    {\bf\Large Universidad de La Frontera}\\[2pt]
    {\small Departamento de Matemática y Estadística}\\[4pt]
    \hrule height 1.5pt
    \vspace{0.3cm}
    {\small\textbf{Asignatura:} Trigonometría y Geometría Analítica \hfill \textbf{Carrera:} Pedagogía en Matemática}
\end{center}

\vspace{0.2cm}
\begin{center}
    {\Large\bfseries Guía de Ejercicios N°1: Puntos y Rectas}
\end{center}
\vspace{0.3cm}

\sffamily\small

\begin{itemize} 
\item[{\textbf A)}] \textbf{Resuelva los siguientes problemas relativos a puntos, distancias y ecuaciones de rectas en el plano:}

\Opensolutionfile{ans}[ans1] 

\begin{ex}
Calcule el perímetro del cuadrilátero $ABCD$ cuyos vértices son $A(6,1)$, $B(3,5)$, $C(-1,-2)$ y $D(2,-4)$. 
\begin{sol} Perímetro = $5+\sqrt{65}+\sqrt{13}+\sqrt{41}$ unidades. \end{sol}
\end{ex}

\begin{ex}
La abscisa de un punto $A$ es $-6$ y su distancia al punto $B(1,3)$ es $\sqrt{74}$. ¿Cuáles son las coordenadas del punto $A$? 
\begin{sol} $A(-6,8)$ o $A(-6,-2)$. \end{sol}
\end{ex}

\begin{ex}
Determine analíticamente si los puntos $A(12,1)$, $B(-3,-2)$ y $C(2,-1)$ son colineales. 
\begin{sol} Son colineales, pues $\overline{AC} + \overline{CB} = \overline{AB}$. \end{sol}
\end{ex}

\begin{ex}
Halle el punto $P$ del eje $X$ que equidista de $Q(5,2)$ y $R(4,-1)$. 
\begin{sol} $P(6,0)$. \end{sol}
\end{ex}

\begin{ex}
Sea el triángulo de vértices $A(2,6)$, $B(-3,2)$ y $C(0,2)$. Calcule el área y el perímetro del triángulo y determine las coordenadas de los puntos medios $E$ de $\overline{CA}$, $F$ de $\overline{BA}$ y $G$ de $\overline{BC}$.
\begin{sol} Área = $6\,\text{u}^2$, Perímetro = $\sqrt{41}+2\sqrt{5}+3$, Puntos Medios: $E(1,4)$, $F\left(-\frac{1}{2},4\right)$ y $G\left(-\frac{3}{2}, 2\right)$. \end{sol}
\end{ex}

\begin{ex}
La recta de ecuación $y+4=m(x+8)$ pasa por el punto de intersección de las rectas $L_1: 2x+3y+5=0$ y $L_2: 5x-2y-16=0$. Calcular el valor de $m$. 
\begin{sol} $m=\frac{1}{10}$. \end{sol}
\end{ex}

\begin{ex}
Los puntos extremos de un segmento son $P_1(2,4)$ y $P_2(8,-4)$. Halle el punto $P(x,y)$ que divide a este segmento en dos partes tales que $\frac{P_1P}{PP_2}=2$. 
\begin{sol} $P\left(6,-\frac{4}{3}\right)$. \end{sol}
\end{ex}

\begin{ex}
Calcule el área del triángulo $\triangle ABC$, cuyos vértices son $A(-2,1)$, $B(4,7)$ y $C(6,-3)$. 
\begin{sol} $36\,\text{u}^2$. \end{sol}
\end{ex}

\begin{ex}
Halle los puntos de la recta $x-y=1$ que están a 2 unidades del punto $(3,0)$.
\begin{sol} $(3,2)$ y $(1,0)$. \end{sol}
\end{ex}

\begin{ex}
Encuentre la ecuación de la recta que pasa por el punto de intersección de las rectas $x-3y+2=0$ y $5x+6y-4=0$ y que es paralela a la recta $4x+y+7=0$. 
\begin{sol} $12x + 3y = 2$. \end{sol}
\end{ex}

\begin{ex}
Halle los valores de $a$ y $b$ para que las ecuaciones $ax+(2-b)y-23=0$ y $(a-1)x+by+15=0$ representen rectas que pasen por el punto $(2,-3)$. 
\begin{sol} $a=4$ y $b=7$. \end{sol}
\end{ex}

\begin{ex}
Halle la distancia del origen a la recta $2x-3y+9=0$. 
\begin{sol} $\frac{9}{\sqrt{13}}$. \end{sol}
\end{ex}

\begin{ex}
Halle dos puntos de la recta $x = y - 5$ cuya distancia a la recta $3x+4y=-6$ sea $3$. 
\begin{sol} $\left(-\frac{41}{7}, -\frac{6}{7}\right)$ y $\left(-\frac{11}{7}, \frac{24}{7}\right)$. \end{sol}
\end{ex}

\begin{ex}
Determine las coordenadas del ortocentro y del baricentro del triángulo cuyos vértices son $A(2,4)$, $B(8,2)$ y $C(4,8)$.
\begin{sol} Ortocentro $H(5,5)$, Baricentro $G\left(\frac{14}{3},\frac{14}{3}\right)$. \end{sol}
\end{ex}

\begin{ex}
Demuestre que el área del triángulo formado por el eje $Y$ y las rectas no paralelas $L_1: y= m_1x + b_1$ y $L_2: y = m_2x + b_2$ está dada por: $A = \dfrac{(b_2-b_1)^2}{2|m_2-m_1|}$.
\begin{sol} Demostración directa evaluando el punto de intersección $\left(\frac{b_1-b_2}{m_2-m_1}, \frac{m_2b_1-m_1b_2}{m_2-m_1}\right)$ y la base en el eje $Y$ de longitud $|b_2-b_1|$. \end{sol}
\end{ex}

\Closesolutionfile{ans}

\vspace{0.3cm}
\item[{\textbf B)}] \textbf{En cada uno de los siguientes casos, halle la ecuación de la recta que satisfaga las condiciones dadas:}

\Opensolutionfile{ans}[ans2] 

\begin{ex}
La pendiente es $4$ y pasa por el punto $(2,-3)$. 
\begin{sol} $y= 4x-11$. \end{sol}
\end{ex}

\begin{ex}
Pasa por los puntos $(3,19)$ y $(-5,4)$. 
\begin{sol} $15x-8y+107=0$. \end{sol}
\end{ex}

\begin{ex}
Pasa por el punto $(1,-7)$ y es paralela al eje $Y$. 
\begin{sol} $x=1$. \end{sol}
\end{ex}

\begin{ex}
Pasa por el punto $(-3,-4)$ y es paralela al eje $X$. 
\begin{sol} $y=-4$. \end{sol}
\end{ex}

\begin{ex}
Pasa por el punto $(1,4)$ y es paralela a la recta $2x-5y=-7$. 
\begin{sol} $2x - 5y = -18$. \end{sol}
\end{ex}

\begin{ex}
Pasa por el punto $(1,-3)$ y es perpendicular a la recta $2y-3x-4=0$. 
\begin{sol} $2x + 3y + 7 = 0$. \end{sol}
\end{ex}

\begin{ex}
Tiene interceptos sobre los ejes $X$ e $Y$ en $2$ y $-3$ respectivamente. 
\begin{sol} $3x - 2y = 6$. \end{sol}
\end{ex}

\begin{ex}
Es simetral (mediatriz) del segmento cuyos extremos son $A(-1,3)$ y $B(2,-3)$. 
\begin{sol} $2x - 4y - 1 = 0$. \end{sol}
\end{ex}

\begin{ex}
Pasa por $P(3,5)$ y la suma de sus interceptos con los ejes coordenados es igual a $12$ (con interceptos no nulos).
\begin{sol} $x + y = 8$ o $5x + 3y = 30$. \end{sol}
\end{ex}

\Closesolutionfile{ans}

\vspace{0.3cm}
\item[{\textbf C)}] \textbf{Determine el valor del parámetro $k \in \mathbb{R}$ para que se cumpla cada condición dada:}

\Opensolutionfile{ans}[ans3] 

\begin{ex}
La recta $kx + (k + 1)y + 3 = 0$ sea perpendicular a la recta $2x - 7y + 2 = 0$. 
\begin{sol} $k = -\frac{7}{5}$. \end{sol}
\end{ex}

\begin{ex}
La recta $4x + 5y + k = 0$ forme con los ejes coordenados un triángulo de área $\dfrac{5}{2}$. 
\begin{sol} $k = 10$ o $k = -10$. \end{sol}
\end{ex}

\begin{ex}
La distancia del punto $(2,-2)$ a la recta $kx + 3y + 3 = 0$ sea igual a $2$. 
\begin{sol} $k = -\frac{9}{4}$ o $k = \frac{7}{12}$. \end{sol}
\end{ex}

\begin{ex}
La recta $kx + (k-1)y - 18 = 0$ sea paralela a la recta $4x + 3y + 7 = 0$. 
\begin{sol} $k = 4$. \end{sol}
\end{ex}

\begin{ex}
La distancia del origen a la recta $x + ky - 7 = 0$ sea igual a $2$. 
\begin{sol} $k = \pm\frac{\sqrt{45}}{2}$. \end{sol}
\end{ex}

\Closesolutionfile{ans}

\end{itemize}

\newpage
\small
\begin{center}
    {\bfseries\Large Soluciones - Guía N°1: Puntos, Rectas y Geometría Analítica en el Plano}\\[10pt]
\end{center}

\textbf{Sección A}
\begin{multicols}{2}
\input{ans1}
\end{multicols}

\vspace{0.4cm}
\textbf{Sección B}
\begin{multicols}{2}
\input{ans2}
\end{multicols}

\vspace{0.4cm}
\textbf{Sección C}
\begin{multicols}{2}
\input{ans3}
\end{multicols}

\end{document}